\subsection*{An explicit temperature bound for tracking}

\propositionheading{Temperature-dependent entropy tracking}{prop:phil-temperature-tracking}
In the four-cycle of \cref{conj:phil-contrast-formation}, let
\(0<\mathrm T\le1/100\), \(p\in(0,1)^4\), and
\(p_{\min}=\min_i\{p_i,1-p_i\}\). Define
\[
 \mathrm K_0=210\sqrt{2\pi/5},\qquad
 \mathrm c_* =\min\left\{
 \frac2{1+4\mathrm K_0e^{13/8}},\frac1{\mathrm K_0e^{13/4}}\right\}.
\]
The entropy production of the full frozen law satisfies
\[
 \mathscr D_p(f)\ge
 \frac{\mathrm c_*\mathrm T^{5/2}}
 {4+3\log(1/\mathrm T)+8\log(1/p_{\min})}
 \operatorname{Ent}_{\mu_p}(f).
\]
Write \(\varepsilon=\epsilon_{\mathrm L}\),
\(t_{\mathrm L}=2\varepsilon^{-1}\log(1/\varepsilon)\), and
\[
 \mathrm D(\mathrm T,\varepsilon)
 =4+3\log(1/\mathrm T)+8\log2+32\log(1/\varepsilon),\qquad
 \ell=\frac{\mathrm c_*\mathrm T^{5/2}}
                {\mathrm D(\mathrm T,\varepsilon)}.
\]
For \(0<\varepsilon<e^{-1}\) and \(0\le t\le2t_{\mathrm L}\),
\begin{equation}
 \operatorname{TV}(q_t,\mu_{p(t)})
 \le e^{-\ell t/2}
       \sqrt{\frac{13(12+|\mathsf m_0|^2)}{4\mathrm T}}
    +\frac{2\varepsilon\mathrm D(\mathrm T,\varepsilon)}
                 {\mathrm c_*\mathrm T^{7/2}}.
\label{eq:phil-temperature-tracking}
\end{equation}
In particular, for \(\varepsilon=\mathrm T^a\) with any \(a>7/2\),
this total variation tends to zero uniformly on
\([t_{\mathrm L},2t_{\mathrm L}]\) and on bounded sets of initial means.
This conclusion compares one-time laws; it supplies no imprint lower bound.
\status{ESTABLISHED}

\paragraph{Proof.}
Use \(u=(\bar k+\bar m)/\sqrt2\),
\(v=(\bar k-\bar m)/\sqrt2\), and \(\omega\), and put
\(R=[D,-\sqrt2I]\). The common kernel
\(K=\ker R\oplus\ker D\) has dimension five. Every
\(\mathsf H_\sigma\) vanishes there, and the reference precision preserves
\(K\) and \(N=K^\perp\). Its restrictions have eigenvalues one and five.
The passive process on \(K\) is therefore an independent Gaussian
Ornstein--Uhlenbeck process with entropy contraction at least \(2\mathrm T\).
The active process on \(N\), together with the assignment, is autonomous:
all redraw costs depend only on its residuals. Rescaling its coordinate by
\(\mathrm T^{-1/2}\) gives unit diffusion and drift matrix
\(a_\sigma=(\mathsf H_\sigma+\mathrm T\mathsf\Omega)|_N\).
Since every edge multiplicity is at most two,
\[
 \mathrm TI\preceq a_\sigma\preceq13I.
\]
Indeed \(RR^\top=DD^\top+2I\) has eigenvalues \(2,4,4,6\),
while \(D^\top D\) has eigenvalues \(0,2,2,4\), so
\(\|\mathsf H_\sigma\|\le12\).

Divide the four determinants in
\cref{eq:phil-selection-sheet-determinants} by the passive determinant
\(5\mathrm T^5\). Their positive coefficients, evaluated at
\(\mathrm T=1/100\) for upper bounds, give the active bounds
\[
 \begin{aligned}
 5120\mathrm T^3&\le d_M\le5262\mathrm T^3,&
 896\mathrm T&\le d_E\le1017\mathrm T,\\
 896\mathrm T&\le d_N\le995\mathrm T,&
 3072&\le d_C<5120.
 \end{aligned}
\]
Thus \(w_M\) and \(w_C\) are the largest and smallest active sheet
weights, and \(w_M\le\mathrm T^{-3/2}/\sqrt{5120}\),
\(w_C/w_M\ge\mathrm T^{3/2}\). The finite-label constant in the
entropy proof consequently satisfies
\[
 L_p\le4+3\log(1/\mathrm T)+8\log(1/p_{\min}).
\]

Apply the Gaussian-sheet and prior-weighted-ball argument of
\cref{eq:phil-selection-gaussian-inequalities,eq:phil-selection-entropy-contraction}
in the seven active coordinates, using the ball of radius \(1/2\), whose
volume is \(\pi^3/840\). Its Gaussian comparison factor is at most
\(4\mathrm K_0e^{13/8}\mathrm T^{-3/2}\), because
\[
 \frac{(2\pi)^{7/2}}{(\pi^3/840)\sqrt{5120}}=\mathrm K_0.
\]
More explicitly, the two energy coefficients in that argument obey
\[
 C_c\le(2+8\mathrm K_0e^{13/8})L_p\mathrm T^{-5/2},\qquad
 C_j\le4\mathrm K_0e^{13/4}L_p\mathrm T^{-5/2}.
\]
The continuous production equals four times its quadratic form and the
jump production bounds four times its quadratic form. Hence the active
rate is at least \(\mathrm c_*\mathrm T^{5/2}/L_p\). Entropy decomposes
over the independent passive and active factors. Cauchy--Schwarz bounds
the energy of a marginal by the corresponding integrated energy, so the
full rate is at least the smaller of the two rates. The displayed active
lower bound is less than \(2\mathrm T\), proving the full-state inequality.
The form-domain extension uses the positive regularization and cutoffs
already used for \cref{eq:phil-selection-entropy-contraction}.

A one-receiver flip joins only the class pairs \(M\)--\(E\),
\(N\)--\(E\), and \(C\)--\(E\). The determinant bounds above make every
ratio of such adjacent weights at most \(1/\mathrm T\).
Writing the evidence as \(p_iZ_i^++(1-p_i)Z_i^-\), the ratio
\(Z_i^+/Z_i^-\) is between the extremal adjacent ratios. Therefore
\(|\partial_{p_i}\log Z|\le1/\mathrm T\), and
\(\|\mathfrak g\|_1\le4/\mathrm T\).
The row lower bound gives the rate floor \(\ell\) on the stated horizon.
The initial entropy is at most
\(13(12+|\mathsf m_0|^2)/(2\mathrm T)\): in the Gaussian density ratio,
\(\mathsf B_\sigma\preceq13I/\mathrm T\),
\(\mathsf Z\le1\), and the remaining quadratic term has nonpositive
expectation. Insert these bounds in
\cref{eq:phil-selection-tracking} and use Pinsker's inequality to obtain
\cref{eq:phil-temperature-tracking}.
For \(\varepsilon=\mathrm T^a\), its second term is
\(O(\mathrm T^{a-7/2}\log(1/\mathrm T))\), while
\(\ell t_{\mathrm L}\) grows as a positive constant times
\(\mathrm T^{5/2-a}\). The exponential therefore dominates the initial
\(\mathrm T^{-1/2}\) factor, proving the last assertion. \(\square\)